% ---------------------------------------------------------------------------------------------------------
% 関連研究
% ---------------------------------------------------------------------------------------------------------
\section{関連研究}
\label{sec:related}
\noindent\textbf{ぶら下がり移動（brachiation）の力学と制御．}\quad
テナガザルの腕渡りは点質量振り子で近似でき，離手・捕捉時の衝突損失が移動コストを決める~\cite{bertram1999,usherwood2003}．
Gomes と Ruina は 5 リンク平面猿モデルで，衝突を避ける「エネルギー零コスト」の腕渡り軌道が存在することを示した~\cite{gomes2005}．
ロボット側では 2 リンク腕渡りロボットのふりこ励振と離手制御~\cite{saito1994,nakanishi2000}が古典で，
近年は連続的な把持・離手を含む全身軌道最適化と学習が試みられている．
本課題は，これらと異なり (i)~両把持点が周期運動し，(ii)~把持容量が 3\,cm 突起で有限，
(iii)~背面（体を 180 度回して）捕捉する，という三つの制約が同時に効く点に特徴がある．
\noindent\textbf{空中でのひねりと姿勢制御．}\quad
角運動量保存下での空中ひねりは，落下猫~\cite{kane1969}以来，体節の相対運動で全身の向きを変える幾何学的位相の問題として
理解されている~\cite{frohlich1979,yeadon1993}．本稿の 3 次元縮約モデルは体幹 6 自由度と左右独立の肩・肘・股・膝をもち，
離手時の角運動量と空中の四肢運動でひねりを生成する．離手・捕捉の左右時間差を許す点で，両手同時の brachiation モデルと異なる．
\noindent\textbf{接触を含む軌道最適化．}\quad
多相の直接コロケーション~\cite{hargraves1987,betts2010,kelly2017}と接触陰的最適化~\cite{posa2014,mordatch2012}は，
接触の生成・消滅を含む全身運動の生成に広く用いられる．本稿は相の順序が既知（振り→飛行→捕捉→保持）なので多相定式化を採り，
Hermite--Simpson の暗黙力学，把持容量のエピグラフ変数，順応捕捉相と離手時刻誤差のロバスト・シナリオ，
多スタートと連続法による下包絡線を組み合わせた．KKT 乗数を包絡線定理~\cite{fiacco1976,buskens2001}で
能力・身体パラメータの感度に翻訳する点が，通常の運動合成と異なる．
\noindent\textbf{運動技能の学習と世界モデル．}\quad
モデルフリー強化学習~\cite{haarnoja2018}や運動模倣~\cite{peng2018}は疎報酬の離散的成功課題で停滞しやすく，
世界モデル~\cite{ha2018,hafner2020,hafner2025,hansen2022}は状態遷移を学んで想像上の経験を増やす．
軌道最適化を教師とする方法には，最適制御解へ方策を回帰させる誘導方策探索~\cite{levine2013}，
訪問状態を専門家がラベル付けする DAgger~\cite{ross2011}，MPC の価値を学習する POLO~\cite{lowrey2019}がある．
これらはいずれも主変数（行動・状態遷移・価値）を回帰する．
価値の勾配を学習信号に使う研究として，値勾配学習~\cite{fairbank2012}，Sobolev 学習~\cite{czarnecki2017}，
確率的値勾配~\cite{heess2015}があり，Pontryagin の最小原理を学習に組み込む研究として
Pontryagin 微分可能計画法~\cite{jin2020}や到達可能性の深層学習~\cite{bansal2021}がある．
制約付き強化学習~\cite{achiam2017,tessler2019}はラグランジュ乗数を報酬の重みとしてオンラインで学習する．
本稿の双対場学習はこれらの延長にあるが，
(i)~学習対象を「厳密オラクルの KKT 双対変数（costate と制約価格）の場」と定め，
(ii)~行動を方策学習なしに Hamiltonian の点ごとの最小化から導き，
(iii)~包絡線定理により身体パラメータへの転移則を場の一部として得る，
という三点を一体の原理として提示し，同一環境上でオラクル・主変数模倣・強化学習と比較する点が新しい．
\noindent\textbf{把持容量の生理．}\quad
3\,cm 前後の縁でのぶら下がりは指屈筋群の等尺性最大筋力で決まり，縁の奥行きと把持形態（ハーフクリンプ／オープン）で
最大指力が大きく変わる~\cite{watts2004,amca2012}．本稿は両手容量を 1300\,N（体重比で約 2\,BW）を基準に，
能力スケールと体格の関数として掃引した．
